\documentclass[10pt,twoside,twocolumn,a4paper]{article}

%% Default action is for double-blinded peer review where 
%% authors, address and acknowledgements sections are hidden 
%\usepackage[]{bpasts}
%% For accepted papers enable option "accepted" to deanonymize
%% your manuscript 
\usepackage[accepted]{bpasts}

\usepackage{t1enc}
\usepackage[utf8]{inputenc}
\usepackage{amsmath,amsfonts}

\usepackage{graphicx}
\usepackage{flushend}
\usepackage{float}
%\usepackage{tabularx}
\usepackage{tabulary}
\usepackage{booktabs} % To thicken table lines
\usepackage{fixltx2e}

\usepackage[colorlinks=true, allcolors=bpastsblue, pdfborder={0 0 0}]{hyperref}

%%%%%%%%%%%%%%%%%%%%%%%%%%%%%%%%%%%%%%%%%%%%%%%%%%%%%%%
%% HEADER OF THE ARTICLE (Change all fields)
\abtitle{P-wave velocity and compressive strength of soil samples for SCA Biorefinery Östrand AB, Timrå}
\title{P-wave velocity and compressive strength of soil samples for SCA Biorefinery Östrand AB, Timrå}

\abauthor{P. Lindh, P. Lemenkova}
\author{Per Lindh$^{1}$$^{2}$\email{per.a.lindh@gmail.com} \\ Polina Lemenkova$^{3}$\email{polina.lemenkova@ulb.be}}

\Address{$^1$ Swedish Transport Administration, Gibraltargatan 7, Malmö, Sweden \\ $^2$ Lund University, Division of Building Materials, Box 118, SE- 221-00, Lund, Sweden \\ $^3$ Université Libre de Bruxelles. École polytechnique de Bruxelles (Brussels Faculty of Engineering), Laboratory of Image Synthesis and Analysis. Campus de Solbosch, Av. Franklin Roosevelt 50, Brussels 1000, Belgium}

\Abstract{On behalf of SCA, SGI has performed geotechnical and environmental geotechnical laboratory tests on dredged material from SCA Östrand pulp mill. The purpose was to investigate the possibilities of using stabilization / solidification of dredged masses and to recommend a combination of binder components with recipe selection. The experimental design is based on a systematic approach of measuring compressive strength of the samples and P-wave velocity in order to evaluate the properties of soils for construction of new biorefinery. Experiments have been carried out on a laboratory scale to show that it is possible to achieve the desired properties with regard to strength and environmental impact through stabilization and solidification of contaminated sediments. Laboratory tests containing geotechnical and environmental tests were performed in two different phases and included p-wave velocity, compressive strength. Quality requirements to be achieved in the field consisted of the following conditions: compressive strength of at least 140 kPa (cu $\geq$ 70 kPa) after 90 days; replacements in the stabilized material \textless 5 cm two years after stabilization / solidification. The experiments were performed on a sediment sample from Sk\"{o}nviken. In light of SGI's experience from other projects, a combination of cement and granulated blast furnace slag was chosen for the tests as a binder. The quality objective regarding permeability constitutes both a geotechnical and environmental technical quality requirement. The experiments showed that an increased water ratio gives a lower strength while maintaining the amount of binder.}

\Keywords{Unconfined compressive strength; Elastic wave; P-wave velocity; stabilized sediments; Solidified soil; Cement; Slag; Applied Geophysics, Civil Engineering; Construction; Stiffness}

%%%%%%%%%%%%%%%%%%%%%%%%%%%%%%%%%%%%%%%%%%%%%%%%%%%%%%%
%% EDITOR SECTION (Do not change)
\vol{XX} \no{Y} \year{2022}
\setcounter{page}{1}
\doi{10.24425/bpasts.2022.DOI}
\begin{document}
\maketitle

%%%%%%%%%%%%%%%%%%%%%%%%%%%%%%%%%%%%%%%%%%%%%%%%%%%%%%%
%% BODY OF THE ARTICLE

\section{Introduction}

Accurate evaluation of the soil strength plays a major role in the design of a safe construction on a foundation ground. The SCA Biorefinery Östrand AB plans to build a new biorefinery for the production of fuel in form of petrol and diesel from the biological raw materials. This construction is planned to extend the existing infrastructure of the pulp mill in Östrand's industrial area located in Timrå municipality, northern Sweden. The soils in the area prepared for the planned biorefinery need to be stabilized and solidified before the construction works. In addition, extra area needs to be cleaned in a coastal area of the Gulf of Bothnia, Baltic Sea, in order to gain more space for the planned new biorefinery, Fig. \ref{fig:01}. 

Assessment of compressive strength of soils is important before the construction works, for evaluating ground stability aimed to ensure the safety of structures. Choosing the most optimal combination of binder (cement, slag or other materials) is the most important parameter of the soil stabilization and solidification, as reported in many previous studies \cite{Lindh2001,Xizhong,LindhWinter}. At the same time, it is important in geotechnical engineering to select the best proportion of binders to ensure the strength of soils. The impact of various combination of binders and ratio of water for stabilising soils can be assessed by measuring P-waves of soil samples comparatively assessed after a curing period. However, when stabilising soil by various ratio of binders, it does not guarantee the best combination of binders based on the theoretical assessment. Instead, cement-slag-water ratio for stabilising soil samples should be evaluated practically, as a series of the \emph{in-situ} tests, in order to ensure the safety of soils before the planned construction works.

\begin{figure}[H]
\begin{center}
\includegraphics[width=0.99\linewidth]{img/Fig_01} 
\end{center}
\caption{Figure showing reinforcement barrier for water and land area. Figure source: SCA Biorefinery Östrand AB}
\label{fig:01}
\end{figure}

Various practical methods such as experimental and empirical \emph{in-situ} testing of soil samples, and computer-based modelling have been employed to analyse the compressive strength of soils. These include estimating the porosity and density of soils \cite{Lemenkov2021a,Uyanik,Lemenkov2021b,Consoli}, comparison of binder combinations \cite{Bernal}, adding alkali additives as cement replacement \cite{Eyo}, economic assessment of soil stabilization including costs and environmental impact \cite{Rocha}. For instance, previous studies \cite{Emeryetal} experimentally examined the suction behavior of soil stabilized with cement and lime through the development of soil-water characteristic curves, to assess bearing capacity of foundations. Microstructure can be considered to estimate the strength development in cement-stabilized silty clay, to prove that cement stabilization improves the soil structure by increasing cementation bonding and reducing the pore space \cite{Horpibulsuk}. Another example \cite{Lindh2021a} aimed to assess binder components (cement, cement kiln dust and slag) for optimal soil stabilisation using marine sediments from the port of Gothenburg. 

Stabilisation of soil by the cementitious binder increases its strength, because large pores in soil material become filled with binder \cite{Lindh202102}. As a result, total pore volume in a soil mass decreases, which makes it suitable for the infrastructure construction: roads, buildings or industrial spaces. Experimental changing of binders produces favourable improvements of soil properties through a combination of various binders (lime, cement, and lime-cement), which increases strength and resistance of soils. For example, \cite{Mahmood} proved that adding lime to the soil mass before stabilising it with cement improves its physical and mechanical properties. Another study \cite{Suetal2009} introduced a solidification/stabilization (S/S) technology for treating soil contaminated with phenol choosing ordinary Portland cement and carbon in various proportions of binder. These tests demonstrated that soil could be effectively treated by the cement/activated-carbon S/S approach. Diamantis et al. \cite{Diamantis}, based on the evaluation of the uniaxial compressive strength, investigated the relationships among the physical, dynamic and mechanical properties of the selected rocks. 

Soil evaluation prior to engineering construction works  is commonly used in civil engineering. It is generally based on the assessment of the empirical relations between soil parameters and geotechnical applications \cite{BrunarskiDohojda,CzarneckiGemert,Kallenetal2016}. It forms the fundamental basis for evaluation of the foundation capacity for estimating its suitability for building structures. The soil may be processed from the in-situ collected samples using laboratory standard tests. However, recently measurements are aimed at the non-destructive and approaches of soil testing. Besides, there are discontinuities in properties of clayey sediments collected from the coastal areas. As an alternative to the traditional methods, the use of geophysical approaches in civil engineering, such as measuring P-waves velocities, becomes a solution to this challenge, because remote sensing methods present more robust and reliable methods in civil engineering. As a response to these needs, many studies \cite{Frattaetal2005,Zhaoetal2022,Zhangetal2021,Guetal2021,Weietal2022}, have applied geophysical methods in geotechnical measurements aimed to improve the evaluation of the soil properties by measuring P-waves propagation velocity in samples.

As a part of the reinforcement work in the coastal area around Östrand, stabilization and solidification of soils were undertaken. The aim of this work was to replace the contaminated masses collected in the seabed as dredged sediments that cannot be dumped into the sea for the environmental reasons. Besides that, the use of binder aimed to increase the compressive strength of the soil masses, reduce sediments, bind the impurities and thereby prevent leaching of heavy metals from the raw soil. By stabilising added dredged masses on site, contaminated masses under the seabed can also be stabilized at the same time.

\section{Methods}

\newpage
\section{Results}

\subsection{P-wave speed and compressive strength phase 1}

The p-wave velocity has been measured on three occasions for all samples and in connection with the uniaxial pressure tests for the relevant samples. The results from the P-wave measurements are reported in Figure 14 and Figure 15. The measurements show small mutual differences within each recipe but a significant difference between the different recipes. For samples in the $L_WH_B$ series, the compressive strength varied between 246 and 2262 kPa. In the $H_WL_B$ series, the compressive strength varied between 96 and 1090 kPa, see Figure 16. The results show high repeatability between experiments even when determining compressive strength. The results from the uniaxial pressure tests are reported in Table \ref{tab:1}. 

\begin{table}
\centering
\caption{Results of the tests on the uniaxial compressive strength}
\label{tab:1}
\setlength\tabcolsep{0pt} 
\smallskip 
\begin{tabular*}{\columnwidth}{@{\extracolsep{\fill}}ccc}
\toprule
Sample & Compressive strength (kPa) & Shear strength (kPa) \\
\toprule
$H_WL_B\_c/s\_70/30\_01$ & 96 & 48 \\ 
\hline
$H_WL_B\_c/s\_50/50\_01$ & 294 & 147 \\
\hline
$H_WL_B\_c/s\_30/70\_01$ & 1046 & 523 \\
\hline
$H_WL_B\_c/s\_70/30\_02$ & 96 & 48 \\
\hline
$H_WL_B\_c/s\_50/50\_02$ & 294 & 147 \\
\hline
$H_WL_B\_c/s\_30/70\_02$ & 1090 & 545 \\
\hline
$L_WH_B\_c/s\_70/30\_01$ & 246 & 123 \\
\hline
$L_WH_B\_c/s\_50/50\_01$ & 718 & 359 \\
\hline
$L_WH_B\_c/s\_30/70\_01$ & 2262 & 1131 \\
\hline
$L_WH_B\_c/s\_70/30\_02$ & 264 & 132 \\
\hline
$L_WH_B\_c/s\_50/50\_02$ & 710 & 355 \\
\hline
$L_WH_B\_c/s\_30/70\_02$ & 2186 & 1093 \\
\midrule
\bottomrule
\end{tabular*}
\end{table}

The results show that a minimum compressive strength of 280 kPa can be contained with both amounts of binder with the recipes c / s 50/50 and c / s 30/70. When using c / s 70/30, 280 kPa in compressive strength is not obtained. 

\subsection{P-wave speed and compressive strength phase 2}

For test series part two, a total of 30 test specimens were manufactured according to Table 2. All the test specimens in phase 2 contained 30\% cement and 70\% slag. 

\begin{table}
\caption{Manufactured specimens in laboratory phase 2. For the $H_WL_B$ combination, no specimens were performed for shaking tests, hence a smaller number of specimens for that combination. }
\label{tab:2}
\smallskip 
\begin{tabular*}{\columnwidth}{@{\extracolsep{\fill}}ccc}
\toprule
Binder / water ratio  & LW (139 \%) & HW (190 \%)  \\
\toprule
$L_B$ (120 kg/m3) & 8 & 6 \\ 
\hline
$EL_B$ (100 kg/m3) & 8 & 8 \\
\midrule
\bottomrule
\end{tabular*}
\end{table}

Figures 17 to 20 show the results from P-wave measurements for the specimens manufactured for phase 2. For the combination $H_WL_B$, see Figure 19, only 6 specimens were manufactured as this combination was also included in phase 1. In phase 2, the specimens show a slightly greater mutual variation. This is because before phase 1, a new mixing routine was introduced with regard to the charging energy. Due to a communication failure, this routine was not applied in phase 2, which resulted in an increased spread within each group. 

\begin{figure}[H]
\begin{center}
\includegraphics[width=0.99\linewidth]{img/Fig_02} 
\end{center}
\caption{P-wave velocity for the combination $L_WH_B$ and the three different recipes, where 1 - 6 is cement / slag 70/30, 7 - 12 is cement / slag 50/50 and 13 - 18 is cement / slag 30/70 . The P-wave velocity is measured on three occasions, 28, 42 and 71 days after mixing}
\label{fig:02}
\end{figure}

\begin{figure}[H]
\begin{center}
\includegraphics[width=0.99\linewidth]{img/Fig_03} 
\end{center}
\caption{P-wave velocity for the combination $H_WL_B$ and the three different recipes, where 1 - 6 is cement / slag 70/30, 7 - 12 is cement / slag 50/50 and 13 - 18 is cement / slag 30/70 . The P-wave velocity is measured on three occasions, 28, 43 and 70 days after mixing.}
\label{fig:03}
\end{figure}

After 85 days of storage, the samples were pressed uniaxially after the P-wave velocity of the respective specimen was determined. The choice of 85 days was due to the holiday closure of the laboratory and was not judged to give any significant difference in the compressive strength. The result is reported in Figure 21. When the data base and that experiments were performed as a complete 22-factor experiment, an analysis of variance (ANOVA) could be performed. The analysis shows that there is only a linear relationship between the amount of binder and the water ratio, see Figure 22.

\begin{figure}[H]
\begin{center}
\includegraphics[width=0.99\linewidth]{img/Fig_04} 
\end{center}
\caption{Compressive strength after 90 days of storage at 20$^{\circ}$C as a function of the P-wave velocity. The reference line is based on the requirement for compressive strength of 140 kPa with a safety factor of 2, i.e. 280 kPa.}
\label{fig:04}
\end{figure}

\begin{figure}[H]
\begin{center}
\includegraphics[width=0.99\linewidth]{img/Fig_05} 
\end{center}
\caption{P-wave velocity of the LWLB combination. The P-wave velocity is measured on two occasions, 28 and 85 days after mixing. }
\label{fig:05}
\end{figure}

\begin{figure}[H]
\begin{center}
\includegraphics[width=0.99\linewidth]{img/Fig_06} 
\end{center}
\caption{P-wave velocity of the $L_WEL_B$ combination. The P-wave velocity is measured on two occasions, 28 and 85 days after mixing. }
\label{fig:06}
\end{figure}

\begin{figure}[H]
\begin{center}
\includegraphics[width=0.99\linewidth]{img/Fig_07} 
\end{center}
\caption{P-wave velocity of the $H_WL_B$ combination. The P-wave velocity is measured on two occasions, 28 and 85 days after mixing}
\label{fig:07}
\end{figure}

\begin{figure}[H]
\begin{center}
\includegraphics[width=0.99\linewidth]{img/Fig_08} 
\end{center}
\caption{P-wave velocity of the $H_WEL_B$ combination. The P-wave velocity is measured on two occasions, 28 and 85 days after mixing}
\label{fig:08}
\end{figure}

\begin{figure}[H]
\begin{center}
\includegraphics[width=0.99\linewidth]{img/Fig_09} 
\end{center}
\caption{UCS after 85 days of storage as a function of P-wave velocity including results from phase 1 The reference line is based on the requirement for compressive strength of 140 kPa with a safety factor of 2, i.e. 280 kPa}
\label{fig:09}
\end{figure}

\begin{figure}[H]
\begin{center}
\includegraphics[width=0.99\linewidth]{img/Fig_10} 
\end{center}
\caption{Evaluation of the factor experiment. The experiment only shows linear effects between water ratio and amount of binder. The result is expected, i.e. an increased water ratio gives a lower strength while maintaining the amount of binder}
\label{fig:10}
\end{figure}

\section{Discussion}

\section{Conclusions}

Various methods of estimating soil parameters for civil engineering have been used successfully for many years and reported in various studies \cite{Wuetal2021,Anagnostopoulos,Grabiec,Lemenkov2021c,Bahar}. However, existing theoretical guidelines and local specifications for the optimum design mixtures in specific soils and solidification process should always be supported by practical assessment of soil stability. This considers regional specifics and properties of local soils, in order to ensure safety of constructions in a given area. This especially concerns coastal areas with clayey, not stable soils, such as coastal areas of the Bothnian Bay, the Baltic Sea. To address this problem, a comprehensive series of laboratory tests has been performed in the Geotechnical Swedish Institute (GSI), to evaluate and establish the optimum binder-water ratio for different soil samples collected from the Östrand's industrial area in Timrå municipality, Sweden.

%%%%%%%%%%%%%%%%%%%%%%%%%%%%%%%%%%%%%%%%%%%%%%%%%%%%%%%
%% ACKNOWLEDGEMENTS
\begin{acknowledgements}

\end{acknowledgements}

%%%%%%%%%%%%%%%%%%%%%%%%%%%%%%%%%%%%%%%%%%%%%%%%%%%%%%%
%% REFERENCES
%% Since January 1st 2021, IEEE Citation Style should be applied.
%\bibliographystyle{IEEEtran}
\bibliographystyle{IEEEtranUrldate}
%\newpage

%% Change to your *.bst file
\bibliography{References}
%\nocite{*}

\end{document}


